% Evaluación Comparativa de Ejecución de Código vs Llamadas Directas
% en Sistemas Agénticos Basados en MCP
% LNCS format - Springer
%
\documentclass[runningheads]{llncs}
%
\usepackage[T1]{fontenc}
\usepackage[spanish]{babel}
\usepackage{graphicx}
\usepackage{booktabs}
\usepackage{amsmath}
\usepackage{multirow}
\usepackage{xcolor}
\usepackage{listings}
\usepackage{hyperref}
\renewcommand\UrlFont{\color{blue}\rmfamily}
\urlstyle{rm}
%
\lstset{
  language=Python,
  basicstyle=\ttfamily\small,
  keywordstyle=\color{blue},
  commentstyle=\color{gray},
  stringstyle=\color{red},
  breaklines=true,
  frame=single,
  numbers=left,
  numberstyle=\tiny\color{gray}
}
%
\begin{document}
%
\title{Evaluaci\'on Comparativa de Ejecuci\'on de C\'odigo vs Llamadas Directas en Sistemas Ag\'enticos Basados en MCP}
%
\titlerunning{Ejecuci\'on de C\'odigo vs Llamadas Directas en MCP}
%
\author{Mauricio D\'iaz Villarreal\inst{1} \and
Jorge Esteban Ram\'irez Shashida\inst{1}}
%
\authorrunning{M. D\'iaz Villarreal, J.E. Ram\'irez Shashida}
%
\institute{Instituto Tecnol\'ogico Aut\'onomo de M\'exico (ITAM),
R\'io Hondo No.~1, Progreso Tizap\'an, 01080 Ciudad de M\'exico, M\'exico\\
\email{mauricio.diaz.villarreal@itam.mx}}
%
\maketitle
%
% ============================================================================
% RESUMEN
% ============================================================================
\begin{abstract}
Los sistemas ag\'enticos basados en modelos de lenguaje grande (LLMs) utilizan
el Model Context Protocol (MCP) para acceder a herramientas externas. Existen
dos paradigmas de ejecuci\'on: llamadas directas a herramientas MCP, donde el
modelo invoca herramientas una por una acumulando resultados en su contexto, y
ejecuci\'on de c\'odigo, donde el modelo genera programas que orquestan m\'ultiples
llamadas en un entorno aislado retornando \'unicamente el resultado final.
Este trabajo presenta una evaluaci\'on comparativa sistem\'atica de ambos enfoques
mediante un dise\~no factorial $2 \times 2$ que combina arquitectura (agente
\'unico vs multi-agente) con modo de ejecuci\'on (MCP directo vs ejecuci\'on de
c\'odigo), evaluado sobre tres tipos de tareas de complejidad creciente ---tarea
simple, encadenada y compleja--- utilizando ocho servidores MCP. Se implementa
OrchAIstra, un framework de evaluaci\'on que mide tokens consumidos, latencia,
precisi\'on y tasa de \'exito mediante calificaci\'on automatizada por LLM.
Los resultados buscan identificar bajo qu\'e condiciones cada paradigma es
\'optimo, proporcionando gu\'ias cuantitativas para desarrolladores de sistemas
ag\'enticos.

\keywords{MCP \and ejecuci\'on de c\'odigo \and tool calling \and sistemas multi-agente \and eficiencia de tokens}
\end{abstract}
%
% ============================================================================
% 1. INTRODUCCI\'ON
% ============================================================================
\section{Introducci\'on}
\label{sec:intro}

Los modelos de lenguaje de gran escala (LLMs) han alcanzado un nivel de
capacidad sin precedentes en los \'ultimos a\~nos. Modelos como GPT-5~\cite{ref_gpt5},
Claude Opus 4.6~\cite{ref_opus} y Gemini 3~\cite{ref_gemini3} demuestran
desempe\~no competitivo en razonamiento matem\'atico, generaci\'on de c\'odigo,
an\'alisis de texto y resoluci\'on de problemas complejos. Sin embargo, conforme
las capacidades fundamentales se consolidan, la frontera de investigaci\'on se
ha desplazado hacia la \textbf{gesti\'on eficiente del contexto}: la ventana de
contexto, el espacio limitado de tokens que un modelo puede procesar en una
sola interacci\'on se ha convertido en el enfoque para el mejoramiento de estos modelos. Cada token consumido representa un costo econ\'omico
directo y una porci\'on del
espacio disponible para informaci\'on relevante.

Este problema se intensifica con la adopci\'on de \textbf{sistemas ag\'enticos},
en los cuales un LLM interact\'ua con herramientas externas como las bases de datos,
APIs, sistemas de archivos, e incluso otros agentes  para completar tareas del mundo real. El mecanismo
dominante para esta interacci\'on es el \textit{tool calling}: el modelo recibe
las definiciones de las herramientas disponibles dentro de su contexto,
selecciona cu\'al invocar, genera los par\'ametros apropiados, recibe el resultado
y contin\'ua con su tarea a resolver. Este proceso consume tokens en cada paso: las
descripciones de herramientas, los par\'ametros serializados, las respuestas
intermedias y el historial acumulado compiten por espacio dentro de la ventana
de contexto.

En 2024, Anthropic introdujo el \textbf{Model Context Protocol} (MCP)~\cite{ref_mcp},
un protocolo abierto que estandariza la comunicaci\'on entre LLMs y servidores de
herramientas mediante una protocolo cliente-servidor basada en JSON-RPC. MCP es el protocolo estandar
que permite a los modelos descubrir e invocar herramientas de manera uniforme,
independientemente del tipo de implementaci\'on, y ha sido adoptado
r\'apidamente por la industria como est\'andar de facto.

Dentro del ecosistema MCP se va a hacer enfoque en dos metodos para la ejecuci\'on de
tareas que requieren herramientas:

\begin{enumerate}
    \item \textbf{Llamadas directas a MCP}: El modelo invoca herramientas
    individuales a trav\'es del protocolo, una por una, acumulando resultados en
    su contexto a lo largo de m\'ultiples turnos de interacci\'on.

    \item \textbf{Ejecuci\'on de c\'odigo}: el modelo genera un programa que
    orquesta m\'ultiples llamadas a herramientas dentro de un entorno de
    ejecuci\'on aislado (\textit{sandbox}), retornando \'unicamente el resultado
    final al contexto del modelo~\cite{ref_codeexec}.
\end{enumerate}

Ambos esquemas resuelven el mismo problema funcional tener un LLM conectado con herramientas externas de manera estandarizada pero difieren radicalmente en c\'omo utilizan la
ventana de contexto. Las llamadas directas cargan todas las definiciones de
herramientas y resultados intermedios en el contexto; la ejecuci\'on de c\'odigo
delega la orquestaci\'on a un programa externo, reduciendo potencialmente el
n\'umero de tokens consumidos a cambio de un overhead de generaci\'on de c\'odigo
y ejecuci\'on en sandbox.

\subsection{Planteamiento del problema}
\label{sec:problem}

El mecanismo convencional de tool calling sobre MCP presenta dos limitaciones
estructurales que comprometen su escalabilidad.

La primera es la \textbf{carga est\'atica de definiciones de herramientas}: todas
las definiciones disponibles ---nombres, descripciones, par\'ametros y schemas---
deben cargarse en la ventana de contexto \textit{antes} de que el modelo procese
el prompt del usuario. En escenarios con decenas o cientos de herramientas, el
modelo procesa miles de tokens de definiciones antes de leer la tarea, la
mayor\'ia de las cuales no ser\'an utilizadas~\cite{ref_codeexec}.

La segunda es el \textbf{tr\'ansito obligatorio de resultados intermedios por el
contexto}: en secuencias donde el resultado de una herramienta alimenta la
entrada de la siguiente, cada resultado intermedio pasa por la ventana de
contexto para ser interpretado y reformulado. Los mismos datos se procesan
m\'ultiples veces, y en el caso de datos extensos, este patr\'on puede exceder
los l\'imites de contexto.

Estos dos problemas comparten una ra\'iz com\'un: el \textbf{acoplamiento entre
la orquestaci\'on de herramientas y la ventana de contexto del modelo}. En el
paradigma convencional, el contexto cumple una doble funci\'on: es
simult\'aneamente el espacio de razonamiento del LLM y el canal de comunicaci\'on
entre herramientas.

Anthropic propone como alternativa la ejecuci\'on de c\'odigo, reportando
reducciones de hasta 98.7\% en consumo de tokens (de 150,000 a 2,000) en
escenarios con grandes conjuntos de herramientas~\cite{ref_codeexec}. Sin
embargo, esta cifra proviene de un caso de uso espec\'ifico y no constituye una
evaluaci\'on sistem\'atica a trav\'es de diferentes tipos de tareas y arquitecturas
ag\'enticas.

La pregunta central de esta investigaci\'on es: \textbf{?`bajo qu\'e condiciones la
ejecuci\'on de c\'odigo a trav\'es de MCP supera a las llamadas directas en
t\'erminos de consumo de tokens, latencia y precisi\'on en sistemas ag\'enticos?}

\subsection{Motivaci\'on}
\label{sec:motivation}

La resoluci\'on de este problema tiene consecuencias tangibles en m\'ultiples
dimensiones. Desde la perspectiva econ\'omica, los proveedores de modelos
facturan por mill\'on de tokens: la diferencia entre una arquitectura que consume
150,000 tokens y una que consume 2,000 por interacci\'on representa \'ordenes de
magnitud en costos operativos a escala. Desde la perspectiva del desarrollador,
el ecosistema carece de criterios emp\'iricos para elegir entre ambos paradigmas;
la decisi\'on arquitect\'onica se basa actualmente en intuici\'on o en las
recomendaciones de un \'unico art\'iculo t\'ecnico~\cite{ref_codeexec}. Desde la
perspectiva acad\'emica, si bien existen benchmarks de tool-use y evaluaciones de
sistemas multi-agente, ninguno se enfoca en el eje de \textit{c\'omo} se ejecutan
las herramientas y su impacto en eficiencia. Esta ausencia de evidencia emp\'irica
es el vac\'io que este trabajo busca llenar.

\subsection{Contribuciones}
\label{sec:contributions}

Este trabajo presenta las siguientes contribuciones:

\begin{enumerate}
    \item Un \textbf{dise\~no experimental factorial} $2 \times 2$ que combina
    arquitectura (agente \'unico vs multi-agente) con modo de ejecuci\'on (MCP
    directo vs ejecuci\'on de c\'odigo), evaluado sobre tres tipos de tareas de
    complejidad creciente.

    \item \textbf{OrchAIstra}, un framework de evaluaci\'on de c\'odigo abierto
    que implementa cuatro tipos de agentes, ocho servidores MCP duales
    (versi\'on directa y versi\'on de ejecuci\'on de c\'odigo), y un sistema de
    calificaci\'on automatizada por LLM.

    \item Un \textbf{an\'alisis cuantitativo} de consumo de tokens, latencia,
    precisi\'on y modos de fallo a trav\'es de las doce combinaciones
    experimentales.

    \item \textbf{Gu\'ias basadas en evidencia} para la selecci\'on del modo de
    ejecuci\'on \'optimo seg\'un las caracter\'isticas de la tarea.
\end{enumerate}

\subsection{Estructura del art\'iculo}
\label{sec:structure}

El resto del art\'iculo se organiza como sigue. La Secci\'on~\ref{sec:background}
presenta los antecedentes sobre MCP, ejecuci\'on de c\'odigo y arquitecturas
multi-agente. La Secci\'on~\ref{sec:design} detalla el dise\~no experimental,
incluyendo la matriz factorial, los servidores MCP y las m\'etricas. La
Secci\'on~\ref{sec:results} reporta los resultados experimentales. La
Secci\'on~\ref{sec:discussion} discute las implicaciones y limitaciones.
Finalmente, la Secci\'on~\ref{sec:conclusions} presenta las conclusiones y
trabajo futuro.

%
% ============================================================================
% 2. ANTECEDENTES Y TRABAJO RELACIONADO (placeholder)
% ============================================================================
\section{Antecedentes y Trabajo Relacionado}
\label{sec:background}

% TODO: Write 2nd - See section_guide.md for detailed content plan

%
% ============================================================================
% 3. DISE\~NO EXPERIMENTAL (placeholder)
% ============================================================================
\section{Dise\~no Experimental}
\label{sec:design}

% TODO: Write 1st (most concrete) - See section_guide.md for detailed content plan

%
% ============================================================================
% 4. RESULTADOS (placeholder)
% ============================================================================
\section{Resultados}
\label{sec:results}

% TODO: After running experiments

%
% ============================================================================
% 5. DISCUSI\'ON (placeholder)
% ============================================================================
\section{Discusi\'on}
\label{sec:discussion}

% TODO: After results

%
% ============================================================================
% 6. CONCLUSIONES (placeholder)
% ============================================================================
\section{Conclusiones}
\label{sec:conclusions}

% TODO: Write last

%
% ============================================================================
% CREDITS
% ============================================================================
\begin{credits}
\subsubsection{\ackname}
Los autores agradecen al Instituto Tecnol\'ogico Aut\'onomo de M\'exico (ITAM)
por el apoyo brindado para la realizaci\'on de este trabajo.

\subsubsection{\discintname}
Los autores declaran no tener conflictos de inter\'es relevantes al contenido
de este art\'iculo.
\end{credits}
%
% ============================================================================
% REFERENCES
% ============================================================================
\bibliographystyle{splncs04}
%
\begin{thebibliography}{20}

\bibitem{ref_gpt5}
OpenAI: GPT-5 Technical Report (2025)

\bibitem{ref_opus}
Anthropic: Claude Opus 4.6 Model Card (2025)

\bibitem{ref_gemini3}
Google DeepMind: Gemini 3 Technical Report (2025)

\bibitem{ref_mcp}
Anthropic: Model Context Protocol Specification.
\url{https://modelcontextprotocol.io/} (2024)

\bibitem{ref_mcpspec}
Anthropic: MCP Python SDK.
\url{https://github.com/modelcontextprotocol/python-sdk} (2024)

\bibitem{ref_codeexec}
Anthropic: Code execution with MCP. Anthropic Engineering Blog (2025)

\bibitem{ref_strands}
AWS: Strands Agents -- Code Interpreter Tool.
\url{https://github.com/strands-agents/tools} (2025)

\bibitem{ref_docker}
Docker Inc.: Docker Engine Security.
\url{https://docs.docker.com/engine/security/} (2024)

\bibitem{ref_e2b}
E2B: Cloud Runtime for AI Agents.
\url{https://github.com/e2b-dev/e2b} (2024)

% TODO: Add as sections are written:
% - AutoGen, CrewAI (multi-agent frameworks)
% - Tool-use benchmarks (BFCL, API-Bank, ToolBench)
% - Token efficiency studies

\end{thebibliography}
%
\end{document}
